% ----------------------------------------------------------------------
% Section 7 — Evaluation
% Rationale: Testability and cost are assessment results rather than
% design decisions; grouping them in a dedicated Evaluation section
% separates "what was built" from "how it was assessed".
% ----------------------------------------------------------------------
\section{Evaluation}\label{sec:evaluation}

\subsection{Testability}\label{sec:testability}
% TODO: remove section information
% Discuss per-service testability: Lambda unit tests (SAM CLI), local
% emulation of DynamoDB / S3 / API Gateway, limitations for IoT Core
% (Device Advisor only), and the role of LocalStack for higher-level
% integration tests.

\subsubsection{Sensor Node Firmware}
The testability of the sensor node, built around the ESP32 microcontroller
(see \cref{sec:hardware}), presents specific challenges due to its embedded nature.
Unit testing of the firmware is implemented using the AUnit
framework~\cite{aunit_testing_arduino}, which is structured such that either the test suite or the main
measurement loop executes, preventing interference between testing and regular operation.

Unit tests are implemented for isolated modules such as the sensor readings,
sensor data aggregation, Wi-Fi connection, basic MQTT connection and
configuration settings. However, exhaustively testing the MQTT
publication functionality described in \cref{sec:mqtt_publication} via unit tests is
not feasible without extensive mocking. Since such mocking would deviate significantly
from real-world behavior, and the underlying MQTT library is already well-tested
and established in the community, complex local test implementations for MQTT were omitted.

Instead, the hardware node employs a direct feedback mechanism for connectivity
validation during its real-world measurement workflow. Upon successfully
establishing a connection to the MQTT broker, the node subscribes to the test
topic to verify the connection's health, ensuring the publish functionality
operates correctly in the physical environment.

\subsection{Cost Estimation}\label{sec:cost_estimation}
% TODO: remove section information
% Present the monthly operating-cost estimate per device under realistic
% assumptions (users, active hours, message volume). Compare Free Tier
% and Post Free Tier costs across all involved AWS services.
% Template available at end of document: tab:aws-costs.

\subsubsection{Hardware Cost}
The one-time hardware cost per sensor node is based on the components used for prototyping.
A breakdown of the individual component costs is presented in \cref{tab:hw-costs}.

\begin{table}[htb]
	\centering
	\caption{Hardware Cost per Sensor Node}
	\label{tab:hw-costs}
	\small
	\begin{tabularx}{0.7\linewidth}{>{\raggedright\arraybackslash}X r r}
		\toprule
		\textbf{Component}         & \textbf{Unit Cost}     \\
		\midrule
		ESP32-DevKitC              & 8\,€                   \\
		Bosch BME680               & 18\,€                  \\
		TSL2561 Light Sensor       & 5\,€                   \\
		DHT11 Sensor               & 3\,€                   \\
		Joy-it KY-026 Flame Sensor & 3\,€                   \\
		Joy-it KY-028 Thermistor   & 3\,€                   \\
		Prototyping equipment      & \(\sim\)3\,€           \\
		\midrule
		\textbf{Total}             & \textbf{\(\sim\)43\,€} \\
		\bottomrule
	\end{tabularx}
\end{table}

